% !TEX program = xelatex
% ============================================================
%  高中研究性学习论文
%  文本分类任务中线性分类器的基线作用研究
%  编译方式: xelatex main.tex （需两次以生成正确的交叉引用）
% ============================================================
\documentclass[11pt, a4paper]{ctexart}

% ---------- 页面设置 ----------
\usepackage[top=2.1cm, bottom=2.1cm, left=2.2cm, right=2.2cm]{geometry}
\usepackage{amsmath, amssymb}
\usepackage{booktabs}
\usepackage{graphicx}
\usepackage{pgfplots}
\pgfplotsset{compat=1.17}
\usepackage{caption}
\captionsetup{font=small, labelfont=bf, skip=4pt}
\usepackage{enumitem}
\setlist{nosep, leftmargin=1.8em}
\usepackage[numbers, sort&compress]{natbib}
\usepackage[hidelinks]{hyperref}
\usepackage{titlesec}
\titlespacing*{\section}{0pt}{1.0ex plus .2ex}{0.6ex plus .2ex}
\titlespacing*{\subsection}{0pt}{0.8ex plus .2ex}{0.5ex plus .2ex}

% 行距略微收紧，保证 3~5 页篇幅
\linespread{1.02}
\setlength{\abovedisplayskip}{3pt}
\setlength{\belowdisplayskip}{3pt}
\setlength{\abovedisplayshortskip}{2pt}
\setlength{\belowdisplayshortskip}{2pt}

% ---------- 标题信息 ----------
\title{\vspace{-1.2cm}\bfseries 文本分类任务中线性分类器的基线作用研究\\
  \large\mdseries On the Role of Linear Classifiers as Baselines in Text Classification}
\author{姓名：邹贝宁\footnote{高中研究性学习论文。实验代码与数据可在项目目录 \texttt{experiment/} 中复现。}\quad
  指导教师：俞金炎}
\date{\today}

\begin{document}
\maketitle
\vspace{-0.8cm}

% ============================================================
%  摘要
% ============================================================
\begin{abstract}
\noindent
近年来，以 BERT 为代表的预训练语言模型几乎成为文本分类任务的标准方案，
而结构简单、训练迅速的线性分类器（如 TF-IDF 特征上的逻辑回归与线性支持向量机）
常常被忽视甚至完全省略。本文以“线性分类器是否仍然是一个有效的基线”为研究问题，
在 4 个公开数据集（AG News、20 Newsgroups、TNews、ChnSentiCorp）上，
系统比较了 2 种线性模型与 2 种神经网络模型（TextCNN、BiLSTM）的分类精度、
训练时间与模型规模。实验结果表明：在中文短文本与二分类情感任务上，
线性模型与神经网络的宏平均 F1 差距通常不超过 1—2 个百分点，
但训练时间往往快 1—2 个数量级，参数量少 2 个数量级以上；
在类别数较多、文本较长的任务上，神经网络才体现出较明显的优势。
本文认为，在报告任何复杂模型的性能之前，都应当先给出一个调好的线性基线，
否则无法判断模型复杂度带来的收益是否真实存在。
\end{abstract}

\vspace{0.3em}
\noindent\textbf{关键词：}\ 文本分类；线性分类器；基线；TF-IDF；逻辑回归

\vspace{0.6em}
\noindent\textbf{Abstract.}\ 
Pre-trained language models such as BERT have become the default choice for text
classification, while simple and fast linear classifiers---logistic regression and
linear SVM on TF-IDF features---are often skipped. This paper asks whether such
linear models still serve as valid baselines. On four public datasets
(AG News, 20 Newsgroups, TNews, ChnSentiCorp) we compare two linear models with two
neural models (TextCNN, BiLSTM) in terms of accuracy, training time and model size.
We find that on Chinese short-text and binary sentiment tasks the macro-F1 gap
between linear and neural models is usually within 1--2 points, whereas linear models
train one to two orders of magnitude faster with two orders of magnitude fewer
parameters; neural models show a clear advantage only on tasks with many classes and
longer documents. We conclude that a well-tuned linear baseline should be reported
before any claim about the benefit of model complexity.

\vspace{0.4em}
\noindent\textbf{Keywords:}\ text classification; linear classifier; baseline; TF-IDF; logistic regression

\vspace{0.8em}
\hrule
\vspace{0.4em}

% ============================================================
\section{引言}
% ============================================================

文本分类是自然语言处理中最基础、也是应用最广泛的任务之一，
其目标是把一段文本划分到预先定义好的类别中，
典型应用包括新闻主题分类、情感分析、垃圾邮件识别等。

2018 年以来，以 BERT~\citep{devlin2019bert} 为代表的预训练语言模型
凭借大规模语料预训练与注意力机制~\citep{vaswani2017attention}，
在各类文本分类基准上不断刷新最好成绩。
受此影响，很多初学者和竞赛参与者会直接跳过传统方法，
把“使用预训练模型”当作解决问题的默认起点。

然而，复杂模型并不总是必要的。
\citet{lin2023linear} 在 ACL 2023 的一篇短文中指出：
线性分类器是文本分类中“常被遗忘的基线”，
在很多数据集上，经过合理调参的线性模型与 BERT 的差距远小于人们预期，
而一些工作中的 BERT 之所以表现不佳，恰恰是因为没有与之对比的基线来暴露配置问题。
\citet{sculley2018winner} 也强调，缺少严格基线会使研究结论充满偏倚。

高中阶段的计算资源与时间都相当有限，
因此“是否真的需要复杂模型”这一问题对我们的实际研究尤为关键。
本文尝试回答以下三个问题：
\begin{enumerate}[label=(\arabic*)]
  \item 在中英文公开数据集上，线性分类器与轻量神经网络在分类精度上差距有多大？
  \item 在训练时间与模型规模上，两类方法的差距有多大？
  \item 在什么条件下应当优先使用线性模型，什么条件下才需要使用神经网络？
\end{enumerate}

本文的主要工作可以概括为三点：
（1）在 4 个公开数据集上复现了“线性基线”这一研究思路，
覆盖英文与中文、二分类与多分类、短文本与长文本；
（2）同时从精度和效率两个维度进行对比，
而不只关注分类准确率；
（3）给出了在有限算力条件下选择模型的可操作建议。

% ============================================================
\section{相关工作}
% ============================================================

\subsection{基于统计特征的线性分类器}

深度学习方法普及之前，文本分类的主流做法是“特征工程 + 线性分类器”。
\citet{manning2008introduction} 系统总结了向量空间模型：
先把文本表示为词袋向量，再用 TF-IDF 加权削弱高频通用词的影响，
最后用线性分类器完成分类。
\citet{joachims1998text} 证明支持向量机在高维稀疏的文本特征上表现优异，
这一结论至今仍然成立。
\citet{joulin2017bag} 提出的 fastText 同样采用“简单线性结构 + 词袋特征”，
以极低计算代价取得了接近深度模型的精度。

\subsection{神经网络文本分类模型与基线问题}

\citet{kim2014convolutional} 提出的 TextCNN 用多个尺寸的一维卷积核
提取局部 n-gram 特征，结构简单、训练快速，是短文本分类的经典强基线；
\citet{hochreiter1997long} 提出的长短期记忆网络（LSTM）通过门控机制
缓解梯度消失，配合双向结构可利用上下文信息。
此后基于 Transformer 的预训练模型~\citep{vaswani2017attention, devlin2019bert}
成为主流，但微调通常需要 GPU 且超参数敏感~\citep{sun2019how}。

\citet{lin2023linear} 指出，线性分类器是文本分类中“常被遗忘的基线”：
在很多数据集上，经过合理调参的线性模型与 BERT 的差距远小于预期，
而缺少基线会使人无法发现复杂模型的配置问题；
\citet{sculley2018winner} 也强调缺少严格基线会使研究结论充满偏倚。
本文把这套思路落在高中可完成的范围内：
用 sklearn 实现线性基线，用 PyTorch 实现两种轻量神经网络，
在统一协议下比较精度与效率。

% ============================================================
\section{方法}
% ============================================================

\subsection{问题定义与 TF-IDF 特征}

给定训练集 $\mathcal{D}=\{(x_i, y_i)\}_{i=1}^{N}$，
其中 $x_i$ 为文本，$y_i \in \{1,\dots,K\}$ 为类别标签，
目标是学到映射 $f: x \mapsto y$ 使其在测试集上准确率尽可能高。
我们把文本切分为词（中文按字符切分）得到词表 $V$，
并用 TF-IDF 把文档 $d$ 表示为向量，其中词 $t$ 的权重为
\begin{equation}
  w_{t,d} = \mathrm{tf}(t,d) \times \mathrm{idf}(t),
  \qquad
  \mathrm{idf}(t) = \ln \frac{1 + |\mathcal{D}|}{1 + |\{d : t \in d\}|} + 1 .
\end{equation}
直观地说，TF-IDF 让“在当前文档中频繁出现、但在整个语料中罕见”的词获得高权重，
从而突出有区分度的关键词。
本文还对词频取亚线性变换（$1+\ln \mathrm{tf}$），
并在特征中同时加入一元词与二元词组。

\subsection{线性分类器}

\paragraph{逻辑回归（LR）}
对 $K$ 类问题，逻辑回归为每类学习一组权重 $\mathbf{w}_k$，
用 softmax 把得分转为概率
$P(y=k\mid\mathbf{x}) \propto \exp(\mathbf{w}_k^{\top}\mathbf{x} + b_k)$，
训练目标是最小化带 $\ell_2$ 正则的交叉熵损失
\begin{equation}
  \mathcal{L}(\mathbf{W}) = -\frac{1}{N}\sum_{i=1}^{N} \ln P(y_i \mid \mathbf{x}_i)
  + \frac{1}{2C}\|\mathbf{W}\|_2^2 ,
\end{equation}
其中 $C$ 是正则化强度的倒数，$C$ 越大模型越倾向于拟合训练数据。

\paragraph{线性支持向量机（SVM）}
线性 SVM 对每类采用“一对多”策略，寻找最大间隔超平面：
\begin{equation}
  \min_{\mathbf{w}_k, b_k}\ \frac{1}{2}\|\mathbf{w}_k\|^2
  + C \sum_{i=1}^{N} \max\!\bigl(0,\ 1 - \tilde{y}_{ik}(\mathbf{w}_k^{\top}\mathbf{x}_i + b_k)\bigr),
\end{equation}
其中 $\tilde{y}_{ik}=+1$ 表示样本 $i$ 属于第 $k$ 类，否则为 $-1$。
合页损失只惩罚间隔不足的样本，因此 SVM 对噪声相对稳健。
两种线性模型的参数量均等于“特征维度 $\times$ 类别数”，
在高维稀疏特征上既高效又有良好的泛化性。

\subsection{对比用的神经网络模型}

\paragraph{TextCNN}
\citet{kim2014convolutional} 先把词映射为 $d$ 维稠密向量，
再用尺寸为 3、4、5 的卷积核各 100 个提取局部短语特征，
对每个通道做全局最大池化后拼接，最后经全连接层分类；
不同尺寸的卷积核相当于不同长度的 n-gram 检测器。

\paragraph{BiLSTM}
BiLSTM 用两个方向的 LSTM~\citep{hochreiter1997long} 分别从前往后、
从后往前读取词序列并拼接隐状态。
本文在其上加入加性注意力池化，使模型能对不同时间步赋予不同权重：
$\alpha_t \propto \exp(\mathbf{u}^{\top}\tanh(\mathbf{W}\mathbf{h}_t))$，
$\mathbf{z} = \sum_t \alpha_t \mathbf{h}_t$。

% ============================================================
\section{实验设置}
% ============================================================

\subsection{数据集}

我们选取 4 个公开数据集（见表~\ref{tab:datasets}），
覆盖英文与中文、二分类与多分类、短文本与长文本两类场景：
AG News~\citep{zhang2015character} 与 20 Newsgroups~\citep{lang1995newsweeder}
为英文新闻类数据，TNews~\citep{xu2020clue} 是中文新闻标题分类任务，
ChnSentiCorp~\citep{tan2008empirical} 是中文情感二分类语料。
所有数据集均按 7:1:2 的比例分层划分为训练集、验证集与测试集，随机种子固定为 42。

% 自动生成，请勿手工修改：数据集统计信息
\begin{table}[htbp]
\centering
\caption{实验数据集统计信息。长度指英文词数或中文字符数。}
\label{tab:datasets}
\small
\begin{tabular}{llccccc}
\toprule
数据集 & 语言 & 训练 & 验证 & 测试 & 类别数 & 平均长度 \\
\midrule
AG News & 英文 & 5473 & 782 & 1564 & 4 & 37.9 \\
20 Newsgroups & 英文 & 1626 & 233 & 465 & 4 & 194.4 \\
TNews & 中文 & 3499 & 500 & 1000 & 10 & 19.9 \\
ChnSentiCorp & 中文 & 5583 & 798 & 1596 & 2 & 95.8 \\
\bottomrule
\end{tabular}
\end{table}


\subsection{评价指标}

我们使用三个指标：宏平均 F1（macro-F1）、训练时间和可训练参数量。
选用宏平均 F1 而非单纯的准确率，是因为它先分别计算每个类别的 F1 再取平均，
在类别不均衡时更能反映模型的真实表现。

\subsection{实现细节}

线性模型使用 scikit-learn~\citep{pedregosa2011scikit} 实现，
TF-IDF 特征维度上限设为 20000，正则化参数 $C$ 在验证集上选取。
神经网络使用 PyTorch~\citep{paszke2019pytorch} 实现，
优化器为 Adam（学习率 $10^{-3}$，权重衰减 $10^{-4}$），
批大小 64，最多训练 8 轮并采用早停（patience = 3）。
神经网络的词表由训练集构建（最小词频 2，上限 30000），
序列长度取训练集长度的 95 分位数。
所有实验在 CPU 上进行，这是高中阶段最现实的实验条件。

% ============================================================
\section{实验结果与分析}
% ============================================================

\subsection{分类精度}

表~\ref{tab:main-results} 与图~\ref{fig:f1} 给出了各模型在 4 个数据集上的宏平均 F1。
\begin{itemize}
  \item \textbf{线性模型在全部 4 个数据集上都不弱于神经网络。}
  其中线性 SVM 在 AG News（91.14）、20 Newsgroups（87.52）与 ChnSentiCorp（88.90）
  上取得了所有模型中的最好成绩；在 TNews 上，
  逻辑回归以 52.99 领先 TextCNN 的 48.82，领先幅度约 4 个百分点。
  \item \textbf{BiLSTM 在 4 个数据集上全部落后于线性基线。}
  在 TNews 上 BiLSTM 仅有 35.77，比线性 SVM 低了约 15 个百分点。
  对于这类短文本（平均 20 字），双向循环结构并没有带来收益，
  反而因为参数量大、训练样本相对不足而更容易过拟合。
  \item \textbf{数据集难度差异显著。}
  TNews 是 10 类中文新闻标题分类，类别之间语义接近（如“文化”与“娱乐”），
  所有模型的 F1 都明显低于其他三个任务，
  说明该任务的困难主要来自数据本身而非模型容量。
\end{itemize}

这一结果与\citet{lin2023linear} 的观察一致：
在中小规模数据集上，只要特征合适、正则化参数合理，
线性分类器完全有能力与结构复杂的神经网络抗衡。

% 自动生成，请勿手工修改：宏平均 F1
\begin{table}[htbp]
\centering
\caption{各模型在不同数据集上的宏平均 F1（\%）。每组最高分以粗体标出。}
\label{tab:main-results}
\small
\begin{tabular}{lcccc}
\toprule
模型 & AG News & 20 Newsgroups & TNews & ChnSentiCorp \\
\midrule
TF-IDF + LR & 90.43 & 85.80 & \textbf{52.99} & 87.71 \\
TF-IDF + SVM & \textbf{91.14} & \textbf{87.52} & 51.14 & \textbf{88.90} \\
TextCNN & 90.45 & 87.32 & 48.82 & 86.21 \\
BiLSTM & 88.63 & 80.69 & 35.77 & 82.26 \\
\bottomrule
\end{tabular}
\end{table}


\begin{figure}[htbp]
\centering
% 自动生成，请勿手工修改：宏平均 F1 对比
\begin{tikzpicture}
\begin{axis}[
  ybar, bar width=8pt, width=\linewidth, height=4.3cm,
  ymin=0, ymax=105, ylabel={宏平均 F1 (\%)},
  symbolic x coords={AGNews,20NG,TNews,ChnSenti},
  xtick=data, x tick label style={font=\small},
  legend style={at={(0.5,-0.20)}, anchor=north, legend columns=4, font=\small, draw=none},
  legend image code/.code={\draw[#1] (0cm,-0.12cm) rectangle (0.30cm,0.12cm);}
  enlarge x limits=0.18, ymajorgrids, grid style={gray!25},
]
\addplot+[fill=black!12, draw=black] coordinates {(AGNews,90.43) (20NG,85.80) (TNews,52.99) (ChnSenti,87.71)};
\addplot+[fill=black!35, draw=black] coordinates {(AGNews,91.14) (20NG,87.52) (TNews,51.14) (ChnSenti,88.90)};
\addplot+[fill=black!60, draw=black] coordinates {(AGNews,90.45) (20NG,87.32) (TNews,48.82) (ChnSenti,86.21)};
\addplot+[fill=black!85, draw=black] coordinates {(AGNews,88.63) (20NG,80.69) (TNews,35.77) (ChnSenti,82.26)};
\legend{线性：LR,线性：SVM,神经：TextCNN,神经：BiLSTM}
\end{axis}
\end{tikzpicture}

\caption{各模型在不同数据集上的宏平均 F1 对比}
\label{fig:f1}
\end{figure}

\subsection{效率对比}

表~\ref{tab:efficiency} 与图~\ref{fig:time} 对比了训练时间与模型规模，
从中可以看到两类方法在代价上的量级差异。

\begin{itemize}
  \item \textbf{训练时间相差 2 个数量级左右。}
  线性 SVM 在 4 个数据集上的训练时间都在 0.02--0.06 秒之间，
  而 TextCNN 需要 4.0--42.7 秒，BiLSTM 需要 7.6--89.3 秒。
  以 20 Newsgroups 为例，线性 SVM 用时 0.02 秒便取得了 87.52 的 F1，
  而 BiLSTM 训练 89.3 秒只得到 80.69。
  \item \textbf{参数量相差 1--2 个数量级。}
  线性模型的参数量约为 $2\times10^4$ 至 $1.0\times10^5$，
  神经网络则为 $4.6\times10^5$ 至 $2.6\times10^6$。
  \item \textbf{在 CPU 环境下这一差距尤其关键。}
  本文全部实验在 CPU 上完成：线性模型的总耗时不到 1 秒，
  神经网络则需要数十秒到一分半。
  对算力有限的高中生或小型项目而言，这意味着能否快速迭代实验。
\end{itemize}

% 自动生成，请勿手工修改：训练时间与参数量
\begin{table}[htbp]
\centering
\caption{训练时间（秒）与可训练参数量。线性模型在该规模数据上的训练时间通常在 0.1 秒量级。}
\label{tab:efficiency}
\small
\begin{tabular}{llcccc}
\toprule
模型 & 指标 & AG News & 20 Newsgroups & TNews & ChnSentiCorp \\
\midrule
TF-IDF + LR & 训练时间 & 0.28 & 0.20 & 0.30 & 0.16 \\
 & 参数量 & 80.0K & 80.0K & 104.3K & 20.0K \\
\midrule
TF-IDF + SVM & 训练时间 & 0.03 & 0.02 & 0.04 & 0.04 \\
 & 参数量 & 80.0K & 80.0K & 104.3K & 20.0K \\
\midrule
TextCNN & 训练时间 & 7.70 & 21.7 & 4.72 & 31.4 \\
 & 参数量 & 1.74M & 2.41M & 457.5K & 559.5K \\
\midrule
BiLSTM & 训练时间 & 14.1 & 75.2 & 9.23 & 71.6 \\
 & 参数量 & 1.92M & 2.59M & 633.4K & 735.7K \\
\bottomrule
\end{tabular}
\end{table}


\begin{figure}[htbp]
\centering
% 自动生成，请勿手工修改：训练时间对比（对数坐标）
\begin{tikzpicture}
\begin{axis}[
  ybar, bar width=8pt, width=\linewidth, height=4.3cm,
  ymode=log, ymin=0.02, ymax=200,
  ylabel={训练时间 (秒，对数坐标)},
  symbolic x coords={AGNews,20NG,TNews,ChnSenti},
  xtick=data, x tick label style={font=\small},
  legend style={at={(0.5,-0.20)}, anchor=north, legend columns=4, font=\small, draw=none},
  legend image code/.code={\draw[#1] (0cm,-0.12cm) rectangle (0.30cm,0.12cm);}
  enlarge x limits=0.18, ymajorgrids, grid style={gray!25},
]
\addplot+[fill=black!12, draw=black] coordinates {(AGNews,0.279) (20NG,0.202) (TNews,0.304) (ChnSenti,0.155)};
\addplot+[fill=black!35, draw=black] coordinates {(AGNews,0.034) (20NG,0.024) (TNews,0.040) (ChnSenti,0.043)};
\addplot+[fill=black!60, draw=black] coordinates {(AGNews,7.703) (20NG,21.741) (TNews,4.721) (ChnSenti,31.443)};
\addplot+[fill=black!85, draw=black] coordinates {(AGNews,14.104) (20NG,75.221) (TNews,9.230) (ChnSenti,71.594)};
\legend{线性：LR,线性：SVM,神经：TextCNN,神经：BiLSTM}
\end{axis}
\end{tikzpicture}

\caption{各模型训练时间对比（纵轴为对数坐标）}
\label{fig:time}
\end{figure}

\subsection{讨论：什么时候需要神经网络？}

结合上述结果，可以给出以下经验判断。

\paragraph{优先使用线性模型的情形}
文本较短、特征词区分度较高时（如 AG News 的新闻标题摘要、
TNews 的中文标题、ChnSentiCorp 的评论短句），
线性模型已经足够，且几乎不需要调参就能得到接近最优的结果。
此外，当训练样本较少时（如 20 Newsgroups 每类仅约 580 条），
线性模型的泛化反而更好。

\paragraph{可以考虑神经网络的情形}
本次实验中神经网络并未取得优势，但这一结论有其边界。
从原理上看，神经网络的价值在于建模词序与长距离依赖：
当文本较长、且判断需要结合上下文的语法结构而非个别关键词时，
卷积与循环结构才更可能带来收益。
本次 4 个数据集中最长的 20 Newsgroups 平均仅 194 个英文词，
且话题词（如“gun”“hockey”“graphics”）区分度很高，
因此线性模型已能覆盖主要信息。

\paragraph{更重要的结论：先做出基线}
本文最有价值的发现其实来自实验过程中的一次失误。
最初实现中文 TF-IDF 时，我们直接沿用了 sklearn 的默认词切分规则，
而中文没有空格，导致一整句被当作一个 token，
特征维度仅为 287，TNews 上的宏平均 F1 只有 9\%（随机基线为 10\%）。
在排查这一问题的过程中我们意识到：
如果一开始就只跑神经网络（它使用了正确的字符级分词，能得到约 49 的 F1），
这个严重的流程缺陷就会被 BERT 或 TextCNN 的“还不错”的表现掩盖。
\citet{lin2023linear} 强调基线的作用是“合理性检查”，
本文的经历正是这一观点的一个具体例证。

\subsection{局限}

本文的结论存在若干限制。
第一，为保证高中阶段可完成，我们没有微调 BERT 等预训练模型，
因此结论只能说明“线性模型不弱于轻量神经网络”，
不能直接推广为“线性模型不弱于预训练模型”。
第二，受算力限制，神经网络的训练轮次与超参数搜索范围都较有限，
其表现可能未达最优。
第三，每个数据集只使用单一随机种子划分，
未报告多次运行的均值与标准差。

% ============================================================
\section{结论}
% ============================================================

本文围绕“线性分类器是否仍是有效的文本分类基线”这一问题，
在 4 个中英文公开数据集上对比了 2 种线性模型与 2 种神经网络模型。
实验表明，线性 SVM 在全部 4 个数据集上都取得了最高或接近最高的宏平均 F1，
且训练时间比神经网络快约 2 个数量级、参数量少 1--2 个数量级。
在算力受限的条件下，线性分类器应当作为默认的起点。

我们由此得到两点体会。
其一，模型复杂度本身不是目的，应当先获得一个调好的简单基线，
再判断增加复杂度是否真的带来收益。
其二，基线不只是“陪跑”的数字，它能暴露数据处理流程中的错误——
本文中中文分词缺陷的发现，正是得益于线性模型异常低的分数。

后续工作可以沿两个方向展开：
一是加入 BERT 等预训练模型的微调结果，使对比更完整；
二是对每个数据集使用多个随机种子重复实验，
并报告置信区间，使结论更加稳健。

% ============================================================
\bibliographystyle{plainnat}
\bibliography{references}

\end{document}
